\subsection{Tecnologías conocidas por el distribuidor}

Durante el proceso de captación de requisitos descrito en la Sección~\ref{sec:RequirementsCaptureMethodology}, se ha identificado que el distribuidor tiene experiencia previa con herramientas de gestión empresarial como Microsoft Excel, Microsoft Access y Google Sheets. Estas herramientas se utilizan principalmente para la gestión de pedidos, el seguimiento de inventarios y la elaboración de informes comerciales.

Si bien estas soluciones permiten gestionar información de forma flexible, presentan importantes limitaciones desde el punto de vista de los sistemas de información. En particular, no existe integración entre las distintas herramientas utilizadas, lo que provoca redundancias en la gestión de datos, inconsistencias en la información y dificultades para disponer de datos actualizados en tiempo real. Además, estas herramientas no incorporan mecanismos estructurados de control de acceso, trazabilidad de operaciones o gestión de usuarios, lo que limita su capacidad para soportar un sistema de gestión comercial completo.

\subsection{Aplicaciones existentes de la empresa central}

La empresa central dispone de una aplicación móvil oficial denominada \textit{KIN Cosmetics}, disponible para dispositivos móviles, que actúa principalmente como un catálogo digital de productos y como repositorio de contenidos formativos para profesionales del sector \cite{kintrends}. A través de esta aplicación es posible consultar las distintas líneas de producto, acceder a información técnica y visualizar material formativo relacionado con técnicas de peluquería.

No obstante, esta aplicación no ofrece funcionalidades orientadas a la gestión operativa del distribuidor, tales como la gestión de pedidos, el control de entregas o la integración con sistemas administrativos, lo que limita su utilidad en el contexto de la actividad comercial diaria.

Por otro lado, la página web oficial de la empresa permite consultar el catálogo de productos e incluye información detallada sobre sus características \cite{kincosmetics_web}. En este caso, sí se contempla la posibilidad de realizar pedidos, aunque el proceso no se encuentra integrado dentro de un sistema propio, sino que se basa en la redirección a plataformas externas de comercio electrónico.

Por ejemplo, el producto \textit{KINESSENCES Nourish Intense Mask} puede visualizarse en la web oficial \cite{kincosmetics_producto}, desde donde se facilita el acceso a su compra a través de plataformas como Amazon \cite{amazon_producto}. Este enfoque permite la adquisición de productos, pero no proporciona herramientas específicas para la gestión comercial del distribuidor ni para el seguimiento de la relación con los clientes.

\subsection{Aplicaciones de gestión en el sector: Booksy}

Adicionalmente, el distribuidor identificó el uso de la plataforma ``Booksy'', una aplicación orientada a la gestión de reservas y citas en el ámbito de la belleza y el bienestar \cite{booksy}. Esta herramienta permite a los profesionales gestionar su agenda, organizar la información de clientes y facilitar la reserva de servicios a través de internet.

Aunque este tipo de aplicaciones introduce funcionalidades digitales relevantes, como la gestión estructurada de clientes o la automatización de procesos de reserva, no cubre las necesidades específicas del distribuidor analizado. En particular, no contempla aspectos clave como la gestión de pedidos a proveedores, el control logístico de entregas o la integración con sistemas de facturación.

\subsection{Limitaciones de las soluciones existentes}

El análisis de las herramientas disponibles pone de manifiesto que, si bien existen soluciones parciales que cubren determinados aspectos de la actividad (gestión de datos, catálogo de productos o gestión de clientes), ninguna de ellas ofrece una solución integrada que permita abordar de forma conjunta la gestión comercial, logística y administrativa del distribuidor.

Asimismo, estas herramientas no incorporan de manera generalizada funcionalidades clave desde el punto de vista de los sistemas de información, como la gestión de usuarios con distintos roles, la trazabilidad de las operaciones o la centralización estructurada de los datos.

Estas limitaciones justifican la necesidad de una solución específica que permita integrar las distintas dimensiones de la actividad del distribuidor en una única plataforma, proporcionando una base sólida para la gestión de la información y la evolución futura del sistema.
